% TEMPLATE-MILC-v3.tex - plantilla maestra UNICA de las guias MILC v3.
%
% La usan las cuatro familias de guias, cada una con su driver:
%   · Grados 6-11          -> scripts/build-guias-g11.py
%   · Bebras               -> scripts/build-guias-bebras.py
%   · Semillero            -> scripts/build-guias-semillero.py
%   · Territorio interior  -> scripts/build-guias-territorio-interior.py
%
% Antes habia tres copias casi identicas (~400 lineas iguales, 6 distintas):
% cada mejora habia que aplicarla tres veces, y en la practica no se hacia
% (el motor de recursos vivia solo en dos de ellas). Lo que varia entre
% familias esta parametrizado en 6 placeholders que cada driver llena
% (se nombran aqui SIN su sintaxis real: el builder reemplaza por texto plano,
% tambien dentro de los comentarios, y un valor multilinea rompe el preambulo):
%
%   HEADER_IZQ        encabezado izquierdo de pagina
%   HEADER_DER        encabezado derecho de pagina
%   PORTADA_ETIQUETA  rotulo sobre el numero grande (GRADO/MOMENTO/MODULO)
%   PORTADA_SERIE     linea de serie de la portada
%   PORTADA_ID        identificador de la guia en la portada
%   PORTADA_CIRCULO   el circulito de grado (°), o vacio si no aplica
%   PDF_TITULO        titulo en texto plano (sin \textbf ni comillas TeX) para
%                     los metadatos del PDF (/Title); lo produce lib_xelatex.texto_pdf
%
% Composicion (auditoria 2026-09-03): las bandas de estacion piden espacio
% (\Needspace*) y prohiben el salto justo despues (\nopagebreak), asi no quedan
% huerfanas al pie; las cajas largas (softbox, drawbox, infoband, puentebox) son
% `breakable`, asi una caja de media pagina ya no arrastra todo a la siguiente
% dejando la anterior al 40 %. Todo texto sobre mostaza o verde va en milcNegro
% (blanco sobre #E5B400 da 1,93:1 y sobre #58A923 2,95:1; ninguno pasa WCAG AA).
% El resto de claves las documenta content/guias/_SCHEMA.md. El script valida
% que no queden placeholders sin reemplazar antes de escribir el archivo final.

% Etiquetado PDF (latex-lab): /Lang, estructura y texto alternativo de las
% figuras, para lectores de pantalla. Debe ir ANTES de \documentclass.
\DocumentMetadata{lang=es-CO,tagging=on}
\documentclass[11pt,letterpaper]{article}
% headheight=24pt: la cabecera derecha lleva dos lineas (identidad de la guia
% y estacion actual). headsep se acorta en la misma medida para que el bloque
% de texto quede exactamente donde estaba con headheight=12pt.
\usepackage[margin=2.0cm,headheight=24pt,headsep=13pt]{geometry}
\usepackage[table]{xcolor}
\usepackage{fontspec}
\usepackage[spanish]{babel}
\usepackage{tikz}
% Librerías TikZ para los diagramas de las guías (bloque `recursos:`):
%  · babel  → babel en español vuelve activos caracteres como «>», y TikZ los usa
%             en sus opciones (>=stealth, ->). Sin esta librería, cualquier
%             diagrama con flechas revienta con «\language@active@arg> has an
%             extra }». Debe ir ANTES de que se dibuje nada.
%  · positioning        → sintaxis «below left=2pt and 3pt».
%  · decorations.pathreplacing → llaves ({brace}) para señalar rangos.
\usetikzlibrary{babel,positioning,decorations.pathreplacing}
\usepackage{graphicx}
% Circuitos: el trabajo de electrónica se hace hoy con micro:bit y sus sensores
% integrados (MakeCode, sin cableado), así que no se necesitan esquemáticos.
% Si algún día entran montajes con protoboard, basta descomentar esta línea:
% los diagramas .tex/.tikz declarados en `recursos:` ya se incrustan solos.
% \usepackage{circuitikz}
\usepackage[most]{tcolorbox}
\usepackage{tabularx}
\usepackage{array}
\usepackage{fancyhdr}
\usepackage{titlesec}
\usepackage[document]{ragged2e}  % bandera a la derecha en todo el cuerpo
\usepackage{enumitem}
\usepackage{microtype}
\usepackage{etoolbox}
\usepackage{needspace}
% hyperref al final del preambulo: metadatos del PDF (titulo, autor, idioma,
% familia), marcadores por estacion (\pdfbookmark) y enlaces sin recuadro.
\usepackage[hidelinks,
  pdftitle={Decidir con miedo: la certeza no llega, y el miedo exagera},
  pdfauthor={PhD. Álvaro Cárdenas Orozco},
  pdfsubject={Territorio interior · Educación socioemocional · Grado 11° · Momento 3 · MILC v3},
  pdflang={es-CO},
  bookmarksopen=true,bookmarksnumbered=false]{hyperref}

% Helvetica Neue no trae la flecha U+2192, asi que cada «→» escrito literalmente
% en el YAML se compilaba a NADA, en silencio (462 flechas en 61 guias: rutas del
% tipo «paleta → tipografia → lockup» salian rotas). Mapeamos el caracter a su
% version matematica, que si tiene glifo. Asi el autor puede seguir escribiendo
% «→» comodamente en el YAML.
\usepackage{newunicodechar}
\newunicodechar{→}{\ensuremath{\rightarrow}}
% Mismo problema con los simbolos de marca y vinetas: el «✓» de «una palomita ✓»
% salia como cuadrito vacio en el PDF de todas las guias que lo usaban.
\usepackage{pifont}
\newunicodechar{✓}{\ding{51}}
\newunicodechar{✗}{\ding{55}}
\newunicodechar{✔}{\ding{52}}
\newunicodechar{•}{\textbullet}
\newunicodechar{≥}{\ensuremath{\geq}}
\newunicodechar{≤}{\ensuremath{\leq}}

\setmainfont{Helvetica Neue}
\setsansfont{Helvetica Neue}
\setlength{\parindent}{0pt}
\setlength{\parskip}{6pt}
% Sin particion de palabras: en 6.º-7.º una palabra cortada por guion al final
% de linea («Comuni-cacion») frena la lectura mucho mas que una linea corta.
\hyphenpenalty=10000
\exhyphenpenalty=10000
\pretolerance=2000
\tolerance=3000
\setlength{\emergencystretch}{3em}
\linespread{1.10}
\renewcommand{\arraystretch}{1.25}
\setlist{leftmargin=5mm,itemsep=1mm,topsep=1mm}
\newcolumntype{Y}{>{\RaggedRight\arraybackslash}X}

\definecolor{milcMagenta}{HTML}{E6007E}
\definecolor{milcVino}{HTML}{5A0038}
\definecolor{milcTurquesa}{HTML}{0093A5}
\definecolor{milcVerde}{HTML}{58A923}
\definecolor{milcMostaza}{HTML}{E5B400}
\definecolor{milcOcre}{HTML}{B86600}
\definecolor{milcNegro}{HTML}{241816}
\definecolor{milcGris}{HTML}{F7F7F4}
\definecolor{milcLinea}{HTML}{E8E2DD}
\definecolor{milcGradePrimary}{HTML}{0D9488}
\definecolor{milcGradeDark}{HTML}{115E59}
\definecolor{milcGradeSoft}{HTML}{CCFBF1}
\definecolor{milcGradeAccent}{HTML}{CFE7FF}
\definecolor{milcGradeText}{HTML}{FFFFFF}
\color{milcNegro}

\pagestyle{fancy}
\fancyhf{}
% Cabecera de dos lineas: identidad de la guia arriba y, a la derecha y en
% gris, la estacion en curso (\markright desde \milcestacion). Antes de la
% primera estacion la segunda linea va vacia; el \strut mantiene la altura
% para que la linea de arriba no baile entre paginas.
\lhead{\small Territorio interior · Educación socioemocional\\[-1pt]{\footnotesize\strut}}
\rhead{\small Grado 11° · Momento 3 · MILC v3\\[-1pt]{\footnotesize\strut\color{milcNegro!65}\rightmark}}
\cfoot{\small\thepage}
\renewcommand{\headrulewidth}{0pt}

% Sobre mostaza (#E5B400) y verde (#58A923) el texto va en milcNegro; sobre el
% resto de la paleta, en blanco. Se usa dentro de fonttitle/fontupper, donde un
% \color posterior manda sobre coltitle.
\newcommand{\milctintasobre}[1]{%
  \ifboolexpr{test{\ifstrequal{#1}{milcMostaza}} or test{\ifstrequal{#1}{milcVerde}}}%
    {\color{milcNegro}}{\color{white}}}

\newtcolorbox{titlebox}[1]{
  enhanced,colback=#1,colframe=#1,arc=7mm,boxrule=0pt,
  left=7mm,right=7mm,top=3mm,bottom=3mm,
  before skip=8pt,after skip=8pt,
  fontupper=\sffamily\bfseries\Large\color{white}
}

\newtcolorbox{softbox}[3]{
  enhanced,breakable,pad at break=3mm,
  colback=#2,colframe=#1,borderline west={4pt}{0pt}{#1},
  arc=2mm,boxrule=.4pt,left=4mm,right=4mm,top=3mm,bottom=3mm,
  before skip=6pt,after skip=8pt,title={#3},coltitle=white,
  colbacktitle=#1,fonttitle=\sffamily\bfseries\milctintasobre{#1},fontupper=\RaggedRight,
  attach boxed title to top left={xshift=3mm,yshift=-2mm},
  boxed title style={arc=2mm, boxrule=0pt}
}

\newtcolorbox{drawbox}[2]{
  enhanced,breakable,pad at break=4mm,
  colback=white,colframe=#1,arc=4mm,boxrule=1pt,
  left=5mm,right=5mm,top=4mm,bottom=4mm,before skip=8pt,after skip=8pt,
  title={#2},coltitle=white,colbacktitle=#1,fonttitle=\sffamily\bfseries\milctintasobre{#1},
  fontupper=\RaggedRight,
  attach boxed title to top left={xshift=4mm,yshift=-2mm},
  boxed title style={arc=3mm, boxrule=0pt}
}

\newtcolorbox{infoband}[1]{
  enhanced,breakable,pad at break=2mm,colback=#1!8,colframe=#1!50,arc=2mm,boxrule=.5pt,
  left=4mm,right=4mm,top=2mm,bottom=2mm,before skip=4pt,after skip=4pt,
  fontupper=\small\RaggedRight
}

% Puente narrativo entre secciones (contrato editorial Conéctate).
% Da continuidad: "Acabas de X. Ahora vas a Y porque Z."
% Visualmente: banda izquierda en color del grado, texto en itálica pequeña.
\newtcolorbox{puentebox}{
  enhanced,breakable,pad at break=2mm,colback=milcGradeSoft,colframe=milcGradeDark,
  borderline west={3pt}{0pt}{milcGradeDark},
  arc=0mm,boxrule=0pt,left=5mm,right=4mm,top=2.5mm,bottom=2.5mm,
  before skip=4pt,after skip=8pt,
  fontupper=\itshape\small\color{milcNegro}\RaggedRight
}
% La flecha va en modo matematico a proposito: Helvetica Neue NO tiene el glifo
% U+2192, asi que \textrightarrow se compilaba a NADA («Missing character: There
% is no → in font Helvetica Neue») y el puente salia sin flecha. En modo
% matematico la flecha viene de la fuente matematica y si se dibuja.
\newcommand{\puente}[1]{%
  \begin{puentebox}%
    \textcolor{milcGradeDark}{$\rightarrow$}\hspace{2mm}#1%
  \end{puentebox}%
}

\newcommand{\linead}[1]{\noindent\color{milcLinea}\rule{#1}{0.5pt}\color{milcNegro}}
\newcommand{\respuesta}[1]{\par\vspace{2mm}\foreach \n in {1,...,#1}{\noindent\color{milcLinea}\rule{\linewidth}{0.5pt}\par\vspace{5mm}}\color{milcNegro}}
\newcommand{\checkbox}{\textcolor{milcGradeDark}{\fbox{\phantom{x}}}\hspace{5pt}}

% ─── Listas aireadas para las guías socioemocionales ────────────────────
% Componer las verificaciones como «\checkbox texto\\» deja el texto que salta
% de línea debajo del cuadrito y apelmaza el bloque. Estos entornos alinean la
% continuación y airean los ítems. Los usa Territorio Interior; cualquier
% familia puede hacerlo.
\newlist{checklista}{itemize}{1}
\setlist[checklista]{
  label=\textcolor{milcGradeDark}{\fbox{\phantom{x}}},
  leftmargin=9mm, labelsep=3.2mm, itemsep=2.4mm,
  topsep=2.5mm, parsep=0pt, partopsep=0pt, before=\RaggedRight
}
\newlist{pasoslista}{enumerate}{1}
\setlist[pasoslista]{
  label=\textcolor{milcGradeDark}{\sffamily\bfseries\arabic*.},
  leftmargin=8mm, labelsep=3mm, itemsep=2.8mm,
  topsep=1mm, parsep=0pt, partopsep=0pt, before=\RaggedRight
}

\newcommand{\phasecard}[4]{
\begin{tcolorbox}[enhanced,colback=#3,colframe=#2,arc=3mm,boxrule=.5pt,height=3.05cm,valign=center,halign=center,left=1.5mm,right=1.5mm,top=2mm,bottom=2mm]
{\sffamily\bfseries\fontsize{22}{24}\selectfont\textcolor{#2}{#1}}\par
{\sffamily\bfseries\small #4}
\end{tcolorbox}
}

% ── Piezas del rediseno 2026 (legibilidad 6.º-11.º) ───────────────────
% Banda de estacion: reemplaza el titlebox macizo. Titulo a la izquierda,
% dato practico (tiempo/modalidad) a la derecha, en una sola linea.
% Nunca huerfana: pide 9 lineas libres (o salta de pagina rellenando con
% \vfill) y prohibe el salto justo despues de la banda. Nueve y no seis:
% la caja partible que sigue necesita ~80pt (titulo adosado, relleno y dos
% lineas) para arrancar en la misma pagina; con menos, tcolorbox la manda
% entera a la siguiente y la banda se queda sola. El penalty 10000 va
% en `after`, ANTES del espacio posterior: un `after skip` (aunque sea 0pt)
% es un \vskip tras la caja, y ese glue es un punto de corte legal que un
% \nopagebreak escrito despues de \end{tcolorbox} ya no cierra. Cada banda
% deja marcador PDF y actualiza la cabecera derecha.
\newcounter{milcestacion}
\newcommand{\milcestacion}[2]{%
\Needspace*{9\baselineskip}%
\stepcounter{milcestacion}%
\pdfbookmark[1]{#2}{milcestacion\themilcestacion}%
\markright{#2}%
\begin{tcolorbox}[enhanced,colback=#1,colframe=#1,arc=2mm,boxrule=0pt,
  left=5mm,right=5mm,top=2.6mm,bottom=2.6mm,before skip=11pt,
  after={\par\nopagebreak\vskip7pt}]
{\sffamily\bfseries\fontsize{15}{18}\selectfont\milctintasobre{#1}#2}
\end{tcolorbox}}

% Tarjeta de paso numerada: sustituye las tablas «Pilar / Aplicacion», que
% eran formato de consulta docente, no de estudiante. El numero va en un
% circulo sobre el borde para que la secuencia se lea de un vistazo.
\newcommand{\milcpaso}[3]{%
\Needspace{4\baselineskip}%
\begin{tcolorbox}[enhanced,colback=white,colframe=milcLinea,arc=1.5mm,boxrule=.9pt,
  left=14mm,right=4mm,top=3mm,bottom=3mm,before skip=4pt,after skip=4pt,
  fontupper=\RaggedRight,
  overlay={\node[anchor=north west,circle,fill=#2,text=white,inner sep=0pt,
    minimum size=7.4mm,font=\sffamily\bfseries\small]
    at ([xshift=3.6mm,yshift=-3.2mm]frame.north west) {#1};}]
#3
\end{tcolorbox}}

% Casilla marcable de tamano util con lapiz (la anterior era \fbox{\phantom{x}},
% de alto variable segun el texto que la seguia).
\newcommand{\milcmarca}{\framebox[3.8mm]{\rule{0pt}{3.4mm}}}

% Fila marcable de lista suelta (checks de la estacion 1).
\newcommand{\milccheckrow}[1]{%
\par\vspace{1.8mm}\noindent
\makebox[6.4mm][l]{\raisebox{-.55mm}{\textcolor{milcNegro}{\milcmarca}}}%
\parbox[t]{\dimexpr\linewidth-6.4mm\relax}{\RaggedRight #1}\par}

% Celda marcable dentro de la rubrica (tabla de autoevaluacion). Misma
% sangria francesa, pero pensada para vivir dentro de una columna X.
\newcommand{\milcopcion}[1]{%
\makebox[5.2mm][l]{\raisebox{-.55mm}{\textcolor{milcNegro}{\milcmarca}}}%
\parbox[t]{\dimexpr\linewidth-5.2mm\relax}{\RaggedRight #1}}

% ── Recursos visuales de la guía ──────────────────────────────────────
% Declarados en el bloque `recursos:` del YAML y emitidos por el builder.
% \guiaFigura   → imagen rasterizada o PDF (foto, carta celeste, montaje).
% \guiaDiagrama → fuente TikZ/circuitikz (.tex) incrustada como vector:
%                 es la vía para circuitos y esquemáticos editables.
% [#1] = texto alternativo (alt) · #2 = ruta relativa al .tex · #3 = pie
% El alt llega al PDF etiquetado (/Alt) y lo lee el lector de pantalla.
\newcommand{\guiaFigura}[3][]{%
\begin{center}
\includegraphics[width=0.88\linewidth,height=0.38\textheight,keepaspectratio,alt={#1}]{#2}\\[1.5mm]
{\footnotesize\itshape\textcolor{milcNegro!75}{#3}}
\end{center}
}

\newcommand{\guiaDiagrama}[2]{%
\begin{center}
\input{#1}\\[1.5mm]
{\footnotesize\itshape\textcolor{milcNegro!75}{#2}}
\end{center}
}

\begin{document}
\thispagestyle{empty}

% ── Portada ───────────────────────────────────────────────────────────
% Antes: un rectangulo de color a pagina completa. Se veia bien en pantalla,
% pero en la fotocopiadora del colegio se comia el toner y salia gris sucio.
% Ahora la banda ocupa el tercio superior y el resto va en blanco.
\begin{tikzpicture}[remember picture,overlay]
  \path[fill=milcGradePrimary]
    (current page.north west) rectangle ([yshift=-10.6cm]current page.north east);

  \node[anchor=north west,text=milcGradeText] at ([xshift=2.0cm,yshift=-1.55cm]current page.north west)
    {\sffamily\bfseries\fontsize{12}{14}\selectfont MOMENTO};
  \node[anchor=north west,text=milcGradeText,opacity=.85] at ([xshift=2.0cm,yshift=-2.35cm]current page.north west)
    {\sffamily\bfseries\fontsize{8.5}{10}\selectfont TERRITORIO INTERIOR · SOCIOEMOCIONAL};
  \node[anchor=north east,text=milcGradeText,opacity=.85,align=right] at ([xshift=-2.0cm,yshift=-1.55cm]current page.north east)
    {\sffamily\bfseries\fontsize{9}{12}\selectfont I.E. Sor María Juliana\\Cartago, Valle del Cauca};

  \node[anchor=west,text=milcGradeText] (gradonum) at ([xshift=1.85cm,yshift=-7.1cm]current page.north west)
    {\sffamily\bfseries\fontsize{112}{112}\selectfont 3};
  % El circulito de grado (°) solo aplica a las guias de grado («11°»); en
  % Bebras y Semillero el numero grande es un momento/modulo, no un grado.
  % Se posiciona relativo al nodo `gradonum`, no en coordenadas absolutas.
  

  \node[anchor=west,text=milcGradeText,text width=11.4cm,align=left] at ([xshift=1.5cm,yshift=0.5cm]gradonum.east)
    {
     {\sffamily\bfseries\fontsize{9}{11}\selectfont\MakeUppercase{Territorio interior · Socioemocional}}\\[3mm]
     {\sffamily\bfseries\fontsize{21}{25}\selectfont\setlength{\baselineskip}{27pt}Decidir\\[4pt]con miedo\\[4pt]la certeza no llega}\\[3.5mm]
     {\sffamily\bfseries\fontsize{15}{18}\selectfont Grado 11° · Momento 3}
    };
\end{tikzpicture}

\vspace*{9.4cm}
\vfill

{\sffamily\bfseries\fontsize{8}{10}\selectfont\textcolor{milcGradeDark}{LO QUE TE LLEVAS HOY}}

\begin{tcolorbox}[enhanced,colback=white,colframe=milcNegro,arc=2mm,boxrule=1.6pt,
  left=5mm,right=5mm,top=4mm,bottom=4mm,before skip=3pt,after skip=12pt,
  fontupper=\RaggedRight\fontsize{11}{15.5}\selectfont]
Tu decisión trabajada: la que tienes pendiente escrita con lo que sabes y lo que no, el peor escenario con su plan de recuperación, la prueba de los diez años, y el primer paso pequeño con fecha.
\end{tcolorbox}

\begin{tcolorbox}[enhanced,colback=milcGris,colframe=milcGris,arc=2mm,boxrule=0pt,
  left=5mm,right=5mm,top=3.5mm,bottom=3.5mm,after skip=12pt]
{\sffamily\bfseries\fontsize{8}{10}\selectfont\textcolor{milcNegro!55}{ESTA GUÍA ES DE}}\\[3mm]
\begin{tabularx}{\linewidth}{X p{3.2cm} p{3.2cm}}
\linead{\linewidth} & \linead{3.0cm} & \linead{3.0cm}\\[-1mm]
{\fontsize{8.5}{10}\selectfont\textcolor{milcNegro!55}{Nombre completo}} &
{\fontsize{8.5}{10}\selectfont\textcolor{milcNegro!55}{Grupo}} &
{\fontsize{8.5}{10}\selectfont\textcolor{milcNegro!55}{Fecha}}\\
\end{tabularx}
\end{tcolorbox}

\vfill

\noindent\textcolor{milcLinea}{\rule{\linewidth}{1pt}}\\[2mm]
{\sffamily\bfseries\fontsize{9.5}{12}\selectfont Autor: Dr. Álvaro Cárdenas Orozco}\\[1mm]
{\sffamily\fontsize{8.5}{11}\selectfont\textcolor{milcNegro!55}{Metodología MILC · Apertura ancestral · Decidir sin certeza · Triángulo Dussel-Estoicismo-Floridi}}

\newpage

\section*{Apertura: Decidir con miedo: la certeza no llega, y el miedo exagera}

\begin{softbox}{milcGradeDark}{milcGradeSoft}{Saber ancestral}
A finales del siglo XIX, miles de familias entraron a las montañas de lo que hoy son el Quindío y el norte del Valle a hacer una cosa muy concreta: \textbf{abrir trocha}. \emph{Tumbar monte, quemar, sembrar, levantar un rancho ---años de trabajo--- con la idea de que la tierra que uno desmontaba terminaba siendo suya.} \textbf{Y mientras ellos abrían monte, en una notaría estaba pasando otra cosa.} \emph{El \textbf{25 de noviembre de 1884}, mediante la \textbf{escritura pública n.º 693} de la notaría de Manizales, Lisandro Caicedo constituyó la \textbf{Sociedad Anónima Burila}: \textbf{cien accionistas} de reconocida influencia económica y política del Cauca y de Caldas,} \textbf{con un enorme territorio reclamado en papel.} \emph{Era una época de desorden ---tres guerras civiles--- y de reparto de baldíos.} \textbf{Y vino el choque:} la empresa \emph{se enfrentó durante años con los colonos y con las autoridades de \textbf{Salento} y \textbf{Calarcá} por esas mismas tierras}, y hubo memoriales, pleitos y presión sobre quienes ya estaban posesionados. \emph{Fíjate en la lección, que es dura y es exacta:} \textbf{abrir trocha no daba la tierra.} \emph{Aquellas familias decidieron sin ninguna certeza ---y la certeza nunca llegó---. Aun así, entraron: porque la alternativa era no entrar, y esa también era una decisión.}
\end{softbox}

\begin{softbox}{milcTurquesa}{milcGris}{Saber contemporáneo}
¿Y qué tan buenos somos prediciendo cómo nos vamos a sentir si algo sale mal? \textbf{Bastante malos, y el error va siempre en la misma dirección.} La investigación sobre \textbf{predicción afectiva} ---qué tan bien anticipamos nuestras emociones futuras--- describe un fenómeno llamado \textbf{sesgo de impacto}: \emph{la tendencia a \textbf{sobreestimar la intensidad y la duración} de nuestros estados emocionales futuros}. \textbf{Se ha encontrado en poblaciones muy distintas} ---estudiantes universitarios, hinchas de deportes, votantes registrados---; \emph{en el estudio que le dio nombre, los hinchas sobreestimaron no solo cuánto se alegrarían si ganaba su equipo, sino \textbf{cuánto les iba a durar} esa alegría}. \textbf{Y se conocen tres causas, las tres útiles para decidir.} \emph{El \textbf{focalismo}:} al imaginar el futuro nos concentramos en el evento y \textbf{subestimamos todo lo demás} que también va a estar pasando en nuestra vida. \emph{La \textbf{negligencia inmunitaria}:} \textbf{subestimamos nuestra propia capacidad de recuperarnos} de lo adverso. \emph{Y no anticipamos} \textbf{qué tan rápido le encontramos sentido} a lo que nos pasa, lo cual acelera la recuperación. \emph{Traducido:} \textbf{cuando te imaginas fracasando, tu cabeza está exagerando ---y de manera predecible---.}
\end{softbox}

\begin{softbox}{milcMostaza}{milcGris}{Pregunta-puente}
\textit{Los colonos entraron a la montaña \textbf{sin ninguna garantía}, y no la hubo. \textbf{¿Qué decisión llevas meses posponiendo} \ldots \textbf{esperando estar seguro}?}
\end{softbox}

\begin{softbox}{milcVerde}{milcGris}{Saber y saber hacer}
Al terminar podrás: (1) \textbf{entender} que la certeza no llega y que esperarla es, en sí, una decisión; (2) \textbf{corregir} tu pronóstico del peor escenario sabiendo que el sesgo de impacto exagera de manera predecible; (3) \textbf{distinguir} lo reversible de lo irreversible, que es lo que de verdad cambia el riesgo; (4) \textbf{trabajar} una decisión real hasta dejarla en un primer paso pequeño y fechado.
\end{softbox}



\small
\begin{tabularx}{\textwidth}{>{\bfseries}p{3.0cm}Y}
\rowcolor{milcGradePrimary!12}Autor & Dr. Álvaro Cárdenas Orozco\\
DBA & Comprende los fenómenos que afectan a su territorio, regula sus emociones ante ellos y acompaña a otros con empatía (transversal 6°--11°, articulado con la política «Escuela, territorio de vida» del MEN, Resolución 006519 de 2025, y la Ley 1620 de 2013)\\
Referentes & Primera ayuda psicológica (OMS, 2011/2012) · Directrices IASC (2007) · USGS y Servicio Geológico Colombiano · Pueblo nasa y la minga (CRIC) · Dussel · Estoicismo · Floridi\\
Producto final & Tu decisión trabajada: la que tienes pendiente escrita con lo que sabes y lo que no, el peor escenario con su plan de recuperación, la prueba de los diez años, y el primer paso pequeño con fecha.\\
\end{tabularx}
\normalsize

\puente{En el momento 2 miraste con qué vara te mides. \textbf{Hoy toca la consecuencia práctica:} \emph{decidir tu vida cuando nadie te puede decir si estás decidiendo bien}.}

\milcestacion{milcGradeDark}{Ruta MILC: así vamos a aprender}
\begin{tabularx}{\textwidth}{>{\centering\arraybackslash}X>{\centering\arraybackslash}X>{\centering\arraybackslash}X>{\centering\arraybackslash}X}
\phasecard{1}{milcVerde}{milcGris}{Escucha\\[-1mm]\small Observo, escucho y pregunto.} &
\phasecard{2}{milcTurquesa}{milcGris}{Sistematización\\[-1mm]\small Decido sin esperar la certeza.} &
\phasecard{3}{milcMagenta}{milcGris}{Praxis\\[-1mm]\small Construyo y aplico.} &
\phasecard{4}{milcVino}{milcGris}{Evaluación\\[-1mm]\small Reflexión liberadora.}
\end{tabularx}


\puente{Primero, la decisión concreta que tienes encima. \emph{Y qué es exactamente lo que te da miedo, que casi nunca es lo que uno cree.}}

\milcestacion{milcVerde}{Estación 1 de 3 · Escucha}

\begin{softbox}{milcVerde}{milcGris}{Escena para escuchar}
\textbf{Actividad 1 · IDENTIFICA --- Lo que estoy posponiendo.} \textbf{Nadie va a leer esta página.} \textbf{Primera parte.} Escribe \textbf{una decisión que tengas pendiente} y que te dé miedo. \emph{Puede ser grande o mediana:} \textbf{qué vas a estudiar, si te vas o te quedas, si dejas algo, si le dices algo a alguien, si te presentas a algo, si dejas una relación.} \textbf{Segunda parte.} Responde: \emph{¿desde cuándo la tienes pendiente?} \textbf{Y después:} \emph{¿qué estás esperando exactamente para decidir?} \textbf{Sé preciso ---«estar seguro», «que alguien me diga», «que se aclare», «que llegue el momento»---.} \textbf{Tercera parte, la clave.} \emph{¿Qué es lo peor que puede pasar si decides y sale mal?} \textbf{Escríbelo con detalle, no con una palabra:} \emph{no «me va mal», sino qué pasaría, quién se enteraría, qué tendrías que hacer después.} \textbf{Y la última:} \emph{¿qué pasa si no decides nunca?} \emph{Escríbelo también, porque eso también es un escenario.}
\end{softbox}

{\sffamily\bfseries\fontsize{8}{10}\selectfont\textcolor{milcNegro!55}{MARCA CUANDO LO HAYAS HECHO}}
\milccheckrow{Escribiste una decisión concreta y desde cuándo la tienes pendiente.}
\milccheckrow{Escribiste \textbf{qué estás esperando exactamente} para decidir.}
\milccheckrow{Describiste el peor escenario con detalle, y también el de no decidir.}

\begin{infoband}{milcVerde}


\textbf{En tu cuaderno (Actividad 1 · IDENTIFICA).} Título: «Actividad 1. Lo que estoy posponiendo». \textbf{Formato:} la decisión + desde cuándo + qué esperas + el peor escenario detallado + el escenario de no decidir. \textbf{Extensión:} media página. \textbf{Sabes que terminaste cuando} escribiste \textbf{qué pasa si no decides nunca}. Esta página es tuya: \textbf{nadie más tiene que leerla si tú no quieres}.
\end{infoband}

\puente{Ya la tienes. Ahora, por qué la certeza no iba a llegar, cuánto exagera el miedo y qué es de verdad irreversible.}

\milcestacion{milcTurquesa}{Fase 2 · Sistematización: entender la decisión}

\begin{softbox}{milcTurquesa}{milcGris}{Cuatro claves sobre decidir}
Cuatro claves para dejar de esperar una certeza que no viene.
\end{softbox}

\milcpaso{1}{milcTurquesa}{\textbf{La certeza no iba a llegar.} \emph{Los colonos entraron a la montaña sin garantía de nada, y no la hubo:} \textbf{mientras ellos abrían trocha, cien accionistas ya tenían la tierra reclamada por escritura desde 1884.} \emph{Muchos perdieron pleitos que ni sabían que existían.} \textbf{Y aun así, la decisión de entrar no fue absurda:} \emph{la alternativa era quedarse sin nada, con certeza}. \textbf{Esa es la estructura de casi toda decisión importante:} \emph{no se elige entre riesgo y seguridad ---se elige entre dos riesgos distintos---}, y uno de ellos, el de no moverse, viene disfrazado de prudencia porque no exige hacer nada.}
\milcpaso{2}{milcTurquesa}{\textbf{El miedo exagera, y se sabe en qué dirección.} El \textbf{sesgo de impacto} dice que sobreestimamos \textbf{la intensidad y la duración} de lo que vamos a sentir. \emph{Y hay tres razones, las tres corregibles.} \textbf{Focalismo:} al imaginar el fracaso te concentras en él y \textbf{olvidas todo lo demás que también estará pasando} ---tus amigos, tus rutinas, otras cosas que te importan---. \textbf{Negligencia inmunitaria:} \emph{subestimas tu propia capacidad de recuperarte}, aunque ya te hayas recuperado de cosas antes. \textbf{Y no anticipas} \emph{qué tan rápido le vas a encontrar sentido a lo que pase}. \textbf{Conclusión práctica:} \emph{el escenario que escribiste en la Actividad 1 \textbf{está inflado}, y no porque seas exagerado: le pasa a todo el mundo, medido.}}
\milcpaso{3}{milcTurquesa}{\textbf{Casi nada es irreversible, y ahí está el verdadero criterio.} \emph{La pregunta útil no es «¿es riesgoso?» ---casi todo lo es--- sino \textbf{«¿se puede deshacer?»}.} \textbf{Cambiar de carrera se puede; perder un año se recupera; una ciudad se deja; un trabajo se cambia.} \emph{Y hay cosas que no:} \textbf{lo que le hace daño irreparable a alguien, lo que compromete la salud, lo que se firma sin entender, lo que no se puede pagar.} \emph{Con lo reversible, la estrategia correcta es \textbf{probar}} ---la información que buscas casi siempre solo aparece después de empezar---. \textbf{Con lo irreversible, la estrategia correcta es \textbf{la demora de Séneca}:} \emph{esperar, consultar, no firmar hoy.} \textbf{Lo que sale caro es aplicar la cautela de lo irreversible a lo reversible,} \emph{y eso es exactamente lo que hace alguien que lleva un año sin decidir algo que podría probar en un mes.}}
\milcpaso{4}{milcTurquesa}{\textbf{No decidir también decide, y encima decide peor.} \emph{Cuando alguien pospone, siente que está \textbf{manteniendo abiertas} las opciones.} \textbf{Casi nunca es así:} \emph{el tiempo cierra puertas solo ---se pasan las fechas, se llenan los cupos, la otra persona se cansa, la vida empuja hacia el camino por defecto---.} \textbf{Y hay una diferencia grande entre las dos maneras de terminar en el mismo lugar:} \emph{quien decidió quedarse puede explicar por qué; a quien se le pasó el tiempo solo le queda la explicación de que no alcanzó a decidir.} \textbf{Aplazar no es neutral:} \emph{es elegir la opción por defecto, sin haberla examinado, y sin poder defenderla después.}}

\begin{drawbox}{milcTurquesa}{Cómo se trabaja una decisión que da miedo}
\begin{pasoslista}
\item \textbf{Separa lo que sabes de lo que no,} y de lo segundo: \emph{¿qué se puede averiguar y qué no se puede saber nunca?}
\item \textbf{Escribe el peor escenario, y después el plan:} \emph{si pasa eso, ¿qué haría?} \emph{Casi siempre hay respuesta, y escribirla desinfla el miedo.}
\item \textbf{¿Es reversible?} \emph{Si sí, probar. Si no, demorar y consultar.}
\item \textbf{La prueba de los diez años:} \emph{¿de qué me voy a arrepentir más: de haberlo intentado o de no haberlo intentado?}
\item \textbf{Reduce a un paso pequeño} que dependa solo de ti y quepa en esta semana.
\end{pasoslista}
\end{drawbox}

\begin{softbox}{milcMostaza}{milcGris}{Errores comunes y su corrección}
\textbf{Esperar la certeza:} no llega, y esperarla ya es decidir. \textbf{Confundir riesgo con irreversibilidad:} lo que decide la estrategia es si se puede deshacer. \textbf{Creer que posponer conserva opciones:} el tiempo las cierra solo. \textbf{Imaginar el fracaso sin plan:} el escenario sin respuesta es el que paraliza. \textbf{Creer el pronóstico del miedo:} está inflado de forma medida. \textbf{Pedir opinión a todo el mundo:} multiplica varas ajenas y no reduce la incertidumbre. \textbf{Un primer paso que depende de otros:} no es un primer paso.
\end{softbox}

\begin{infoband}{milcTurquesa}


\textbf{En tu cuaderno (Actividad 2 · EXPLICA).} Título: «Actividad 2. El sesgo de impacto». \textbf{Formato:} explica las tres causas ---focalismo, negligencia inmunitaria y la rapidez con que damos sentido--- y \textbf{corrige} con ellas el peor escenario que escribiste en la Actividad 1. \textbf{Extensión:} media página. \textbf{Sabes que terminaste cuando} puedes explicar por qué \textbf{no decidir también es decidir}.
\end{infoband}

\puente{Ya sabes cómo funciona. Es hora de trabajar tu decisión hasta que quepa en un paso que puedas dar esta semana.}

\milcestacion{milcMagenta}{Fase 3 · Praxis: trabajar la decisión}

\begin{softbox}{milcMagenta}{milcGris}{Cómo se trabaja una decisión que da miedo}
Aquí no se busca la decisión correcta ---nadie puede saberla---. Se busca dejarla trabajada hasta que quepa en un paso que puedas dar esta semana.
\end{softbox}

\milcpaso{1}{milcMagenta}{\textbf{Lo que sé y lo que no (8 min).} Haz dos columnas con tu decisión. \emph{Y a la columna de lo que no sabes agrégale una marca:} \textbf{qué se puede averiguar} ---preguntando, buscando, hablando con alguien que ya lo hizo--- \textbf{y qué no se puede saber nunca}. \emph{Casi siempre aparece la sorpresa:} \textbf{una parte grande de lo que te frena \emph{se puede averiguar esta semana} y no lo has hecho.}}
\milcpaso{2}{milcMagenta}{\textbf{El peor escenario, con plan (10 min).} Reescribe el peor escenario \textbf{corregido por el sesgo de impacto} ---acordándote del focalismo y de tu capacidad de recuperarte--- y, al frente, \textbf{qué harías si pasara}. \emph{Concreto:} \textbf{a quién le contarías, qué harías el mes siguiente, cómo lo arreglarías.} \textbf{Este es el ejercicio que más desinfla el miedo,} \emph{porque el miedo vive de escenarios sin respuesta.}}
\milcpaso{3}{milcMagenta}{\textbf{La prueba de los diez años (5 min).} Pregúntate: \emph{a los veintisiete, ¿de qué me voy a arrepentir más: de haberlo intentado o de no haberlo intentado?} \textbf{Escribe la respuesta.} \emph{No es una pregunta retórica ---a veces la respuesta honesta es «de haberlo intentado», y eso también es información valiosa---.}}
\milcpaso{4}{milcMagenta}{\textbf{El paso pequeño (10 min).} Reduce la decisión a \textbf{una acción concreta que puedas hacer esta semana} y que \textbf{dependa solo de ti}: \emph{escribir el correo, hacer la llamada, buscar el requisito, hablar con alguien que ya lo hizo, inscribirte para averiguar, tener la conversación.} \textbf{Con día y hora.} \emph{Y ojo con la trampa más común:} \textbf{«decidir» no es un paso.} \emph{Un paso es algo que otra persona podría ver que hiciste.}}

\begin{drawbox}{milcMagenta}{Antes de cerrar tu decisión}
\begin{checklista}
\item Separaste lo que sabes de lo que no, y marcaste \textbf{qué se puede averiguar}.
\item Reescribiste el peor escenario \textbf{corregido} y con \textbf{plan de respuesta}.
\item Determinaste si la decisión es \textbf{reversible} y ajustaste la estrategia.
\item Respondiste la prueba de los diez años, aunque la respuesta te sorprenda.
\item Tu paso cabe en esta semana y \textbf{depende solo de ti}.
\item Tu paso es algo que \textbf{otra persona podría ver que hiciste}.
\end{checklista}
\end{drawbox}

\begin{softbox}{milcMagenta}{milcGris}{Plantilla de trabajo}
\textbf{Las dos columnas.} \emph{«Sé \_\_\_\_. No sé \_\_\_\_ ---y de eso, puedo averiguar \_\_\_\_---.»} \\[1mm]
\textbf{El peor escenario con plan.} \emph{«Si pasa \_\_\_\_, yo \_\_\_\_.»} \\[1mm]
\textbf{El paso.} \emph{«El \_\_\_\_ a las \_\_\_\_ voy a \_\_\_\_.»} ---algo que se pueda ver que hiciste---.
\end{softbox}

\begin{infoband}{milcMagenta}


\textbf{En tu cuaderno (Actividad 3 · CREA).} Título: «Actividad 3. Mi decisión trabajada». \textbf{Formato:} las dos columnas con lo averiguable + el peor escenario corregido con su plan + si es reversible + la prueba de los diez años + el paso con día y hora. \textbf{Extensión:} una página. \textbf{Sabes que terminaste cuando} tu paso \textbf{tiene fecha}.
\end{infoband}

\puente{El producto termina en un paso pequeño. \emph{Las decisiones grandes no se toman: se van tomando.}}

\milcestacion{milcMostaza}{Producto de la sesión}
\textbf{Decisión trabajada: lo sabido y lo averiguable, el peor escenario con plan, y un primer paso fechado}.
\begin{softbox}{milcMostaza}{milcGris}{Criterios de calidad}
(1) Distingue lo que no se sabe pero se puede averiguar de lo que no se puede saber. (2) El peor escenario está corregido por el sesgo de impacto y tiene un plan de respuesta concreto. (3) Se clasifica la decisión como reversible o irreversible y se ajusta la estrategia. (4) La prueba de los diez años está respondida honestamente, incluso si contradice la intuición inicial. (5) El primer paso es observable, depende solo de quien decide y tiene fecha.
\end{softbox}

\puente{Antes de cerrar, revisa lo que hace fracasar esto: si tu primer paso depende de que antes ocurra algo que no controlas, no es un primer paso.}

\milcestacion{milcVino}{Evaluación liberadora · cinco dimensiones}

\begin{softbox}{milcVino}{milcGris}{Sentido de la evaluación}
Evaluar no es solo revisar una nota. Es reconocer qué aprendí, qué cambió en mí, qué emociones manejé, cómo me relaciono con la ciudad y la comunidad, y qué le dejaré a quien viene detrás.
\end{softbox}

\textbf{Mi reflexión liberadora en cinco dimensiones.} Completa cada frase iniciada:

\small
\begin{tabularx}{\textwidth}{>{\bfseries\RaggedRight\arraybackslash}p{3.4cm}Y>{\RaggedRight\arraybackslash}p{4.6cm}}
\rowcolor{milcVino}\color{white}Dimensión & \color{white}\bfseries Mi respuesta (frame starter) & \color{white}\bfseries Algo que descubrí\\
1. Personal & Lo que más cambió en mí fue \linead{6.5cm} & \linead{4.2cm}\\
2. Emocional & La emoción más fuerte fue \linead{2.5cm} y la regulé \linead{3cm} & \linead{4.2cm}\\
3. Ciudadana & En democracia esto importa porque \linead{6cm} & \linead{4.2cm}\\
4. Local (Cartago) & En mi barrio sería útil para \linead{6.5cm} & \linead{4.2cm}\\
5. Intergeneracional & A mi abuela/abuelo le diría \linead{6.5cm} & \linead{4.2cm}\\
\end{tabularx}
\normalsize

\begin{infoband}{milcVino}
\textbf{Para la próxima sustentación.} Vuelve a esta tabla en seis meses. Verás que las respuestas cambian — no porque lo aprendido cambie, sino porque tú cambias. La evaluación liberadora es una conversación contigo a través del tiempo.
\end{infoband}


\puente{Cerramos con tres voces: el que está más allá de la totalidad, lo que sí está en tu poder, y el cambio permanente de las cosas.}

\milcestacion{milcGradeDark}{Triángulo de pensamiento · cierre filosófico}

\begin{softbox}{milcVino}{milcGris}{Dussel · lente del nosotros}
\textit{«Analéctico quiere indicar el hecho real humano por el que todo hombre, todo grupo o pueblo se sitúa siempre más allá (aná-) del horizonte de la totalidad.»}

{\footnotesize\itshape\color{milcNegro!70}--- Enrique Dussel · Filosofía de la liberación (1977), §2.4}

\emph{La historia de la Burila enseña algo que Dussel formula en abstracto:} \textbf{el horizonte que uno alcanza a ver no es todo lo que hay}. \emph{Los colonos veían monte y trabajo; no veían una escritura firmada en una notaría de Manizales que ya decidía sobre su futuro.} \textbf{Y a los diecisiete pasa algo parecido:} \emph{las opciones que ves son las que alguien te mostró ---las carreras que conoces, los caminos que hay en tu familia, lo que se hace en tu pueblo---,} \textbf{y hay muchas más de las que aparecen en tu horizonte.} \emph{Por eso el paso más rentable de esta guía casi nunca es decidir:} \textbf{es \emph{averiguar}}, \emph{que es la manera práctica de correr el horizonte antes de escoger dentro de él.}

\textbf{Las opciones que estoy considerando, ¿son todas las que hay o solo las que he visto?}
\end{softbox}
\linead{\linewidth}\\[1mm]
\linead{\linewidth}

\begin{softbox}{milcMostaza}{milcGris}{Estoicismo · Epicteto · Enquiridión, 1 (c. 125 d.C.; trad. desde E. Carter) · cuidado interior}
\textit{«Algunas cosas están en nuestro poder y otras no. Están en nuestro poder la opinión, el impulso, el deseo, el rechazo; en una palabra, todo aquello que es obra nuestra.»}

{\footnotesize\itshape\color{milcNegro!70}--- Epicteto · Enquiridión, 1 (c. 125 d.C.; trad. desde E. Carter)}

\emph{Aquí Epicteto resuelve el problema de fondo de esta guía.} \textbf{El resultado de una decisión \emph{no} está en tu poder} ---puedes escoger bien y que salga mal, y escoger mal y que salga bien---. \textbf{Lo que sí es obra tuya es \emph{cómo decides}:} \emph{averiguar lo averiguable, escribir el peor escenario, consultar a quien ya pasó por ahí, no firmar lo que no entiendes, dar el paso.} \emph{Y esa distinción quita un peso enorme,} \textbf{porque el miedo a decidir casi siempre es miedo al \emph{resultado} ---que no controlas---} \emph{y no al proceso, que sí.} \textbf{Una decisión se juzga por cómo se tomó, no por cómo salió.}

\textbf{¿Estoy tratando de garantizar el resultado o de decidir bien?}
\end{softbox}
\linead{\linewidth}\\[1mm]
\linead{\linewidth}

\begin{softbox}{milcTurquesa}{milcGris}{Floridi · lente de la infoesfera}
\textit{«Las tecnologías digitales están re-ontologizando la naturaleza misma de la infoesfera.»}

{\footnotesize\itshape\color{milcNegro!70}--- Luciano Floridi · La cuarta revolución (2014; trad.)}

\textbf{Hoy hay más información disponible que nunca para decidir, y sin embargo cuesta más decidir que antes.} \emph{No es una paradoja:} \textbf{la abundancia produce la ilusión de que en alguna parte está \emph{la} respuesta}, y entonces uno sigue buscando en vez de escoger. \emph{Es una versión elegante de la evitación del momento 1 del grado 10:} \textbf{investigar se siente productivo y aplaza la decisión igual que cualquier otra anestesia.} \emph{Y hay algo más:} \textbf{lo que aparece cuando uno busca «qué estudiar» está optimizado para retener la atención, no para ayudar a decidir.} \emph{Una conversación de veinte minutos con alguien que ya hizo lo que estás pensando hacer vale más que veinte horas de búsqueda,} \textbf{y esa conversación es, casi siempre, el paso pequeño que hay que dar.}

\textbf{¿Estoy investigando para decidir, o investigando para no decidir todavía?}
\end{softbox}
\linead{\linewidth}\\[1mm]
\linead{\linewidth}

\begin{infoband}{milcNegro}
\textbf{Nota para el docente.} Las tres son citas textuales y llevan su referencia debajo. Puedes remitir a la obra en clase.
\end{infoband}

\milcestacion{milcVino}{Mi autoevaluación · marca una casilla por fila}

{\small
\setlength{\extrarowheight}{2.2pt}
\begin{tabularx}{\textwidth}{>{\RaggedRight\arraybackslash\bfseries}p{3.0cm}YYY}
\rowcolor{milcVino}\color{white}\bfseries Criterio & \color{white}\bfseries Lo logré & \color{white}\bfseries En proceso & \color{white}\bfseries Necesito apoyo\\[.6mm]
Comprensión del tema & \milcopcion{Explico las ideas con mis palabras.} & \milcopcion{Reconozco las ideas, falta precisión.} & \milcopcion{Necesito repasar lo básico.}\\[.8mm]
\rowcolor{milcGris}Calidad de la decisión trabajada & \milcopcion{Separo lo que sé de lo que no, tengo plan para el peor escenario y un primer paso con fecha.} & \milcopcion{Analizo la decisión, pero sigo esperando estar seguro para actuar.} & \milcopcion{Sigo posponiendo, y creo que eso es no haber decidido.}\\[.8mm]
Aplicación y producto & \milcopcion{Mi producto cumple los criterios.} & \milcopcion{Está armado, con ajustes pendientes.} & \milcopcion{Aún está en construcción.}\\[.8mm]
\rowcolor{milcGris}Reflexión 5 dimensiones & \milcopcion{Respondí cada una con honestidad.} & \milcopcion{Respondí algunas con detalle.} & \milcopcion{Mis respuestas son muy generales.}\\[.8mm]
Triángulo de pensamiento & \milcopcion{Conecté las 3 voces con mi trabajo.} & \milcopcion{Conecté al menos una.} & \milcopcion{No las relaciono con mi proceso.}\\[.8mm]
\end{tabularx}}

\vspace{3mm}
\textbf{Hoy entendí que:}\\[1mm]
\linead{\linewidth}\\[1mm]
\linead{\linewidth}

\textbf{Mi compromiso para la próxima sesión es:}\\[1mm]
\linead{\linewidth}\\[1mm]
\linead{\linewidth}

\begin{softbox}{milcMostaza}{milcGris}{Cita de despedida · Marco Aurelio}
\textit{``El obstáculo en el camino se vuelve el camino. Lo que estorba al camino se vuelve camino.''}\\[1mm]
Cada error de hoy, bien observado, se convierte en pieza del camino que pisarás mañana.
\end{softbox}

\small
\begin{tabularx}{\textwidth}{>{\bfseries}p{3.6cm}Y}
\rowcolor{milcGradeDark!12}Nombre completo: & \linead{10cm}\\
Fecha: & \linead{4cm} \quad Curso/grupo: \linead{4cm}\\
Mi firma: & \linead{9cm}\\
\end{tabularx}
\normalsize

\begin{infoband}{milcGradeDark}
\textbf{Para la próxima sesión.} Dos cosas. \textbf{Primero, da el paso} ---el que escribiste con fecha--- y trae \textbf{qué averiguaste}. \emph{Casi siempre pasa lo mismo y conviene comprobarlo uno mismo: después de dar el primer paso, la decisión se ve distinta, no porque haya llegado la certeza, sino porque apareció información que no existía mientras estabas quieto.} \textbf{Segundo, pregunta en tu casa por una decisión sin certeza:} averigua con alguien mayor \textbf{una decisión grande que haya tomado sin saber cómo iba a salir} ---venirse a vivir aquí, dejar un trabajo, meterse a estudiar, casarse, separarse, poner un negocio--- y pregúntale \textbf{qué sabía y qué no}, \textbf{qué era lo que más miedo le daba} y \textbf{si eso que temía llegó a pasar}. \textbf{Trae:} lo que averiguaste y lo que te contaron. \emph{Vas a encontrar el sesgo de impacto contado en primera persona: casi nadie recuerda que lo peor haya pasado como lo imaginaba, y casi todos dicen que se habrían decidido antes si hubieran sabido lo rápido que uno se acomoda.}
\end{infoband}

\end{document}
